\chapter{Keselamatan dan prosedur uji terbang}
\begin{bahaya}[Pilot tetap penanggung jawab]
PASUPATASTRA tidak menggantikan pilot. Selama drone dikendalikan dari GCS, pilot memegang remote, melihat drone secara langsung, dan siap mengambil alih kapan saja. Mematikan interlock hanya menolak perintah \emph{baru}; misi yang sedang berjalan dihentikan dari remote (tombol jeda/rem atau RTH sesuai manual vendor) atau tombol \ui{Tahan} di konsol.
\end{bahaya}

\section{Lapisan keselamatan}
\begin{enumerate}
\item \textbf{Interlock di remote}: kendali dari GCS mati secara bawaan.
\item \textbf{Batas misi ganda} di server dan bridge (Bab~4.2).
\item \textbf{Konfirmasi geser} untuk aksi kritis di konsol.
\item \textbf{Log audit}: setiap perintah beserta pengguna, hasil, dan pesan.
\item \textbf{Peran pengguna}: akun pengamat tidak dapat mengirim perintah.
\item \textbf{Failsafe vendor}: hilang sinyal RC memicu kembali ke home (diatur di misi).
\end{enumerate}

\section{Daftar periksa sebelum terbang}
\begin{itemize}
\item Area terbuka tanpa orang di bawah jalur, izin terbang sesuai ketentuan setempat.
\item Pilot memegang remote, satu pengamat visual, komunikasi dengan operator konsol.
\item Baterai drone dan remote di atas 60~\%, GPS minimal 12 satelit, titik home tercatat.
\item Ketinggian RTH di aplikasi vendor di atas rintangan tertinggi; batas ketinggian drone diatur.
\item Konsol: drone \emph{terhubung}, posisi, arah, dan baterai sama dengan aplikasi vendor.
\item Interlock \textbf{mati} sampai pilot menyatakan siap.
\end{itemize}

\section{Urutan uji lapangan pertama}
Lakukan berurutan; lanjut hanya bila tahap sebelumnya berhasil.
\begin{longtable}{@{}p{1.1cm}p{4.2cm}p{9.6cm}@{}}
\toprule
Tahap & Uji & Langkah dan hasil yang diharapkan \\ \midrule
\endhead
0 & Telemetri saja & Interlock mati. Gerakkan drone di darat/terbang manual; konsol mengikuti. Tombol aksi abu-abu dengan alasan \emph{Kendali belum diizinkan}. \\
1 & Arm/disarm (Autel) & Interlock menyala, drone di darat. \ui{Aktifkan motor} → baling-baling berputar; dalam 5 detik \ui{Matikan motor}. Pada DJI, perintah ini ditolak dengan penjelasan (perilaku benar). \\
2 & Lepas landas pendek & \ui{Lepas landas} → drone naik sekitar 1,2 m dan hover. Tunggu 10 detik, \ui{Mendarat}. Periksa kejadian di Log. \\
3 & RTH & Pilot menerbangkan manual ke 15--20 m tinggi dan sekitar 20 m jauhnya. \ui{Kembali ke home} → drone naik ke ketinggian RTH, kembali, dan mendarat. \\
4 & Tahan & Saat RTH berjalan, tekan \ui{Tahan} → RTH dibatalkan, drone hover. \\
5 & Waypoint dekat & Rencana 2--3 titik, ketinggian 20--30 m, jarak kurang dari 100 m, kecepatan 3 m/s, akhir: kembali ke home. Lepas landas dulu, lalu \ui{Unggah ke drone} dan \ui{Mulai misi}. Uji \ui{Jeda} dan lanjutkan. Kejadian \emph{Titik n tercapai} muncul di Log. \\
\bottomrule
\end{longtable}

\section{Jika terjadi masalah}
\begin{itemize}
\item Drone bergerak tidak sesuai: pilot segera ambil alih dengan remote (jeda/rem, lalu stik manual).
\item Konsol kehilangan link: drone tetap mengikuti perintah terakhir dan failsafe vendor; kendalikan dari remote.
\item Catat waktu, lalu ekspor kejadian dari halaman Log untuk evaluasi.
\end{itemize}
